\section{Introduction}
\label{sec:intro}

Let $\{P_t\}_{t \geq 0}$ be a family of probability measures on $\R^d$
evolving over time. The object of interest is the lagged transport geometry
encoded by
\[
    D_j \;:=\; W_2(P_{t-j}, P_t),\qquad j = 1, 2, \ldots, L.
\]
We study the case in which this curve follows the power law
\begin{equation}
    D_j = \zeta\,j^H,\qquad \zeta > 0,\; H \in (0,1),
    \label{eq:power-law}
\end{equation}
and is observed only through noisy estimates $\Dhj$ based on $n$ samples per
lag. In log scale, the observation noise depends on the same family being
inferred, so the problem is heteroscedastic and parameter-coupled.

This coupling produces one information law for the whole paper:
\[
    \mathcal I_{n,L}(H) = n \sum_{j=1}^L j^{2H},
\]
up to scale factors depending on $\zeta$ and $\sigma_0$. The same quantity
governs estimation of $H$, the residual law, the minimax separation
boundary, and adaptivity across lag regimes. Fixed $L$ and growing $L$ are
therefore two regimes of one inferential scale.

The main contributions are:
\begin{enumerate}[leftmargin=*,label=(\arabic*)]
\item \textbf{One feasible inference layer.} The two-step FWLS estimator is
asymptotically oracle-equivalent, yielding a joint inferential law for the
power-law parameters and an exact oracle law for the residual statistic with a
feasible asymptotic counterpart.

\item \textbf{One optimality scale.} The minimax separation boundary is of
order $\bigl(n\sum_{j \leq L} j^{2H}\bigr)^{-1/2}$, with fixed-$L$ and
growing-$L$ rates obtained as special cases.

\item \textbf{One adaptive procedure.} The same inferential scale attains the
boundary without prior knowledge of $H$.

\item \textbf{One unknown-scale extension.} A decoupled Gaussian benchmark for
unknown $\sigma_0$ yields an exact split-$F$ law, while same-sample parametric
bootstrap restores near-nominal calibration on the tested grid.

\item \textbf{Three empirical signatures.} The simulations show the oracle law,
FWLS-oracle agreement, and a boundary transition aligned with the information
scale.
\end{enumerate}

\paragraph{Relation to prior work.}
Direct estimation of roughness exponents for Gaussian processes is classical
\citep{istas1997quadratic,coeurjolly2001estimating}. The present problem differs
because the exponent is inferred from estimated transport discrepancies, so the
noise is itself parameter-coupled. The surrounding transport-rate literature
provides the relevant first layer \citep{fournier2015rate,weed2019sharp,
goldfeld2024limit,delbarrio2023improved}, while locally stationary drift and
adaptive inference connect through \citep{dahlhaus1997fitting,
vogt2012nonparametric,tinio2025bounds,brutsche2024similarity}.
The temporal-validity horizon law of \citet{Parreira2026May} gives one
downstream use of the pair $(H,\zeta)$: those parameters govern the structural
scale of finite memory under drift. The present paper studies the inferential
layer for the lag geometry itself.

Section~\ref{sec:model} defines the observation model and the information
scale. Section~\ref{sec:clt} gives the feasible inference theorem.
Section~\ref{sec:test} establishes the residual law and its response
under departures. Section~\ref{sec:lower} proves the separation lower bound,
Section~\ref{sec:adaptivity} extends it to the growing-lag regime,
Section~\ref{sec:sims} provides the empirical signatures, and
Section~\ref{sec:application} gives the unknown-scale extension.
